\documentclass[aps,reprint]{revtex4-1}
% Engine-specific settings
% Detect pdftex/xetex/luatex, and load appropriate font packages.
% This is inspired by the approach in the iftex package.
% pdftex:
\ifx\pdfmatch\undefined
\else
    \usepackage[T1]{fontenc}
    \usepackage[utf8]{inputenc}
\fi
% xetex:
\ifx\XeTeXinterchartoks\undefined
\else
    \usepackage{fontspec}
    \defaultfontfeatures{Ligatures=TeX}
\fi
% luatex:
\ifx\directlua\undefined
\else
    \usepackage{fontspec}
\fi
% End engine-specific settings
\usepackage[english]{babel}
\usepackage{csquotes}
% \usepackage[backend=biber, sortcites]{biblatex}
\usepackage{url}
\usepackage{textcomp}
\usepackage[usenames,dvipsnames,svgnames, table]{xcolor}
\usepackage[font={scriptsize}]{caption}
\usepackage{amsmath} \usepackage{amsthm} \usepackage{amsfonts}
\usepackage{amssymb}
\usepackage{enumerate}
\usepackage{tikz} \usepackage{float}
\usepackage[procnames]{listings}
\usepackage{pstool} \usepackage{pgfplots}
\usepackage{wrapfig} \usepackage{graphicx} \usepackage{epstopdf}
\usepackage{afterpage}
\usepackage{physics}
\usepackage{multirow}
\usepackage{gensymb}
\usepackage{algorithm}
\usepackage{microtype}
\usepackage[noend]{algpseudocode}
\usepackage{xcolor,colortbl}
\usepackage{microtype}
\usepackage{geometry}
\usepackage{hyperref}
\usepackage{graphicx}
\usepackage{caption}
\usepackage{subcaption}
\usepackage{lipsum}
% \usepackage{pythontex}
% \usepackage{authblk}
\usepackage{nth}
\usepackage{siunitx}
% \usepackage[toc,page]{appendix}
\floatstyle{plaintop}
\restylefloat{table}

% SET LISTING STUFF (FRED)
\definecolor{dkgreen}{rgb}{0,0.6,0}
\definecolor{dred}{rgb}{0.545,0,0}
\definecolor{dblue}{rgb}{0,0,0.545}
\definecolor{lgrey}{rgb}{0.9,0.9,0.9}
\definecolor{gray}{rgb}{0.4,0.4,0.4}
\definecolor{darkblue}{rgb}{0.0,0.0,0.6}
\lstdefinelanguage{cpp}{
      backgroundcolor=\color{lgrey},
      basicstyle=\footnotesize \ttfamily \color{black} \bfseries,
      breakatwhitespace=false,
      breaklines=true,
      captionpos=b,
      commentstyle=\color{dkgreen},
      deletekeywords={...},
      escapeinside={\%*}{*)},
      frame=off,
      language=C++,
      keywordstyle=\color{purple},
      morekeywords={BRIEFDescriptorConfig,string,TiXmlNode,DetectorDescriptorConfigContainer,istringstream,cerr,exit},
      identifierstyle=\color{black},
      stringstyle=\color{blue},
      % numbers=left,
      % numbersep=5pt,
      numberstyle=\tiny\color{black},
      rulecolor=\color{black},
      showspaces=false,
      showstringspaces=false,
      showtabs=false,
      stepnumber=1,
      tabsize=5,
      title=\lstname,
    }

% Custom commands
\newcommand{\unit}[1]{\:\mathrm{#1}}
\newcommand{\noref}[1]{\hyperref[#1]{\ref*{#1}}}
\newcommand{\nonref}[1]{\hyperref[]{\ref*{#1}}}
\newcommand\blankpage{%
  \null
  \thispagestyle{empty}%
  \addtocounter{page}{-1}%
  \newpage}

% Default fixed font does not support bold face
\DeclareFixedFont{\ttb}{T1}{txtt}{bx}{n}{7} % for bold
\DeclareFixedFont{\ttm}{T1}{txtt}{m}{n}{7}  % for normal

\newcommand\numberthis{\addtocounter{equation}{1}\tag{\theequation}}
\DeclareCaptionFont{white}{\color{white}}
\DeclareCaptionFormat{listing}{\colorbox{gray}{\parbox{\columnwidth}{#1#2#3}}}
\captionsetup[lstlisting]{format=listing,labelfont=white,textfont=white}


% Biber for references
% \bibliographystyle{aipauth4-1}

\begin{document}
\sisetup{detect-all}
\title{Direct and iterative approaches to solving Schrödinger's equation for
two electrons in a harmonic oscillator potential}
\author{Erlend Lima}
\author{Frederik J. Mellbye}
\author{Aram Salihi}
\affiliation{University of Oslo, Oslo, Norway \\ Source code available at: \url{https://github.com/Caronthir/FYS3150/tree/master/Project2}}
\date{\today}

\begin{abstract}
  This is really abstract
\end{abstract}
\maketitle
\tableofcontents
\makeatletter
\let\toc@pre\relax
\let\toc@post\relax
\makeatother

\newpage

% % Get all of the data from the script into a dictionary
% \begin{pycode}
% from subprocess import Popen, PIPE
% from latextools import untag_all
% p = Popen(['python', '../analysis/analyze.py', '../cpp'],
% stdout=PIPE, stderr=PIPE)
% output, err = p.communicate()
% data = output.decode('utf8')
% latex_data = untag_all(data)
% \end{pycode}
%
% \begin{table}[ht]
%   \centering
%   \pyc{print(latex_data['error_table'])}
%   \caption{Summary of the errors}
%   \label{tab:error}
% \end{table}
%
% \begin{figure}[ht]
%   \centering
%   \includegraphics[scale=0.5]{figures/function.eps}
%   \caption{\label{fig:general} Caption}
% \end{figure}
\section{Introduction}
\label{sec:introduction}
Eigenvalue problems are among the most occuring equations in
mathematics and science. Many problems are on or can be re-written to the
form
\begin{align*}
  A \mathbf{x} = \lambda \mathbf{x}
\end{align*}
For matrices where $n \geq 5$ there is
no general analytical procedure to obtain the characteristic polynomial, by the
Abel-Ruffini theorem. In most situations the eigenvalue problems have much
larger matrices, and numerical algorithms are needed to solve the problems. In
this paper, the direct Jacobi rotation and iterative Lanczos algorithms
are investigated.

The methods are then used to solve Schrödinger's equation for two electrons in a
three-dimensional harmonic oscillator potential, which can be reformulated into
an eigenvalue problem through scaling and mathematical manipulations. The
energy eigenvalues and associated eigenstates are then obtained by Jacobi's method.

The study of electrons confined in a potential is a highly active research field,
because of the possible nano-medical applications, and in the development of semi-conductors
and quantum computers.
\section{Theory}
\label{sec:theory}
\subsection{Jacobi's rotation algorithm}
A detailed description of Jacobi's rotation algorithm is available at REFERANSE?.
The general idea is rather straightforward, and is based on converting the
matrix $A$ to approximately a diagonal matrix through a series of orthogonal
matrix transformations on the form:
\begin{align*}
  T = S^T A S
\end{align*}
where $S$ is a rotation matrix. The rotation angle is conveniently chosen to
systematically eliminate off-diagonal elements in $A$, until the largest
offdiagonal element is below some threshhold $\epsilon$. Because the rotations
are unitary, the 
\subsection{Idiotisk navn for å understreke at dette er midlertidig}
A unitary transformation preserves orthogonality and the dot product. To see
this, consider an orthogonal basis \(\{\vb{v}_{i}\}\) and an unitary
transformation \(U\). Defining \(\vb{w}_{i}=U\vb{v}_{i}\), the dot product can
be computed as
\begin{align*}
  \vb{w}_{i}\vdot \vb{w}_{j} &= \vb{w}_{i}\vb{w}_{j}^{T}\\
                             &= \left( U\vb{v}_{i} \right)\left( U\vb{v}_{j} \right)^{T}\\
                             &= U\vb{v}_{i}\vb{v}_{j}^{T}U^{T}\\
                             &= U\delta_{ij}U^{T}\\
                             &= \delta_{ij}UU^{T}\\
                             &= \delta_{ij}
\end{align*}
where the last step follows since \(U\) is unitary.
\section{Method}
\label{sec:method}
\section{Results}
\label{sec:results}
\section{Discussion}
\label{sec:discussion}
\section{Conclusion}
\label{sec:conclusion}
\bibliography{references}
\blankpage
\appendix
\section{LOL}
\blankpage
\end{document}
